\chapter*{Introduction}\addcontentsline{toc}{chapter}{Introduction}\label{ch:intro}

For a subgroup $\langle G\rangle$ of prime order $r$ on an elliptic curve
$E/\FF_q$, Pollard rho~\cite{pollard1978,vanoorschot1999} computes discrete
logarithms in about $\sqrt{\pi r/(2A)}$ group operations, where $A$ is the
number of efficiently computable automorphisms~\cite{wiener1998,gallant2000}.
No algorithm is known to beat it on curves of cryptographic shape.

Index calculus breaks finite-field discrete logarithms in subexponential
time. It reached elliptic curves through Semaev's summation
polynomials~\cite{semaev2004} and the extension-field analyses of
Gaudry~\cite{gaudry2009} and Diem~\cite{diem2011}. It is the method that
should exploit the most structure. The question is \emph{where it crosses
rho}: on which curves, at which sizes, under which accounting.

This treatise records what one research programme learned about that
question between July and October 2026. The programme spanned three
repositories (\texttt{crypto}, \texttt{crypto-autoresearcher},
\texttt{cryptanalysis}), several thousand commits and several hundred
preregistered studies. Two companion documents organise the same work
differently:
\begin{itemize}
\item a programme record~\cite{icprogram2026}, chronological and
  exhaustive, with study catalogues;
\item an improvement ledger~\cite{icledger2026}, organised as a ratchet of
  incumbents.
\end{itemize}
This document organises it as a treatise:
\begin{itemize}
\item the theory that constrains index calculus on elliptic curves;
\item the algorithms engineered against those constraints;
\item the measured frontier against a fairly priced rho;
\item what remains open.
\end{itemize}
It runs no new experiment. It re-verified one certificate
(\S\ref{sec:fd-binary}), adds no improvement of its own, and cites the
committed artefact behind every number.

\section*{The argument in brief}

\paragraph{1. Where the cost lives (Part~\ref{part:foundations}).} An index
calculus pays for factor-base set-up, relation collection, linear algebra
and the target's descent. Relation collection costs
(trials per relation) $\times$ (cost of one decomposition). Chapter~\ref{ch:model}
writes this out and maps every lever the programme pulled to the term it
moves.

\paragraph{2. What is fixed for every curve (Chapter~\ref{ch:cover}).} The
Galois group of the summation cover is $(\ZZ/2)^{m-2}$ for every curve,
field and characteristic. Its fibres over $\FF_q$ split by an exact law,
and they split completely on the factor-base locus. The mean decomposition
yield at fixed base size is a double count. Three consequences follow:
\begin{itemize}
\item no curve has anomalous monodromy to exploit;
\item no fibre census can find an anomalous relation rate;
\item no factor-base geometry can raise mean yield.
\end{itemize}
The design space that remains is the cost of the oracle, the algebraic
difficulty of the decomposition system, the linear algebra, and where the
problem lives.

\paragraph{3. How hard is the decomposition system (Chapter~\ref{ch:falldegree})?}
Every subexponential claim for prime-degree binary curves assumes that the
solving degree stays near the first fall degree. In three presentations we
prove that the first fall degree is a constant of the presentation, and
show that the last fall degree, which equals the solving degree for
unsatisfiable systems, is not bounded by it:
\begin{itemize}
\item it is proved to differ for an explicit-membership chain ($(4,5)$
  exactly at $B=3$);
\item it is certified to differ, by independently checked closure
  certificates, for binary Weil descent ($\ge5$ against 2 at $n=33$--$45$);
\item it is measured to differ for prime-field Kummer towers (6 against 4).
\end{itemize}
A proof that the last fall degree grows without bound remains the open
problem, as Huang, Kosters and Yeo~\cite{hky2015} and Kousidis and
Wiemers~\cite{kousidis2019} left it.

\paragraph{4. What engineering bought (Part~\ref{part:algorithms}).}
Root-finding oracles made the decomposition step two to five orders of
magnitude cheaper where they apply. Gröbner engines improved one to two
orders in counted work. Factor bases closed under maps of known logarithm
cut linear algebra by constant factors. On the cold one-target frontier
none of this changed the conclusion: the stages these changes leave alone
(relation yield per probe, table set-up) now dominate.

\paragraph{5. Where index calculus wins (Chapter~\ref{ch:extension} and
Chapter~\ref{sec:transfer}).} Only in weak extension-field families does
the measured cost cross or approach a fairly priced rho:
\begin{itemize}
\item $E_0/\FF_{q^3}$ through a genus-3 GHS cover, by double-large-prime
  hyperelliptic index calculus or by Diem's plane-quartic index calculus;
\item the Joux--Vitse cover of $E(\FF_{p^6})$;
\item possibly $k=4$ over $\FF_{(2^m)^4}$.
\end{itemize}
These are known theoretical results. The contribution here is matched,
priced, verified measurement. None of these families is deployed.

\paragraph{6. What the measurements say (Part~\ref{part:frontier}).} At a
matched reference, rho is cheaper on every prime and prime-degree binary
curve measured: $4.7$--$18\times$ on Koblitz curves of 37--47 bits, and
$10.5$--$102\times$ on prime curves of 20--40 bits, with the gap widening as
$r^{1/6}$. Every apparent online win disappeared against Bernstein--Lange
rho with equal advice~\cite{bl2012}. Matching the reference, pricing native
work and replaying records each moved recorded gaps upward. Chapter~\ref{ch:audit}
lists every such correction.

\section*{Conventions}

\begin{itemize}
\item \emph{Unit.} Costs are group-addition equivalents per $\sqrt r$
  ($S$) unless a family uses its own counted unit, such as $\FF_q$
  multiplications.
\item \emph{Ratios.} ``$X\times$'' is baseline over candidate when it
  describes an improvement, and IC over rho when it describes a distance.
  Each sentence says which.
\item \emph{Intervals.} Intervals are 95\%, as the source recorded them.
\item \emph{Change classes.} Every change is an \cls{advance} (an exponent
  or crossover moves under complete accounting), \cls{engineering} (a
  constant moves), \cls{relabelling} or \cls{accounting}. Against rho, the
  programme produced no \cls{advance} on any prime-order curve of
  cryptographic shape.
\item \emph{Paths.} Artefact paths are relative to the \texttt{crypto}
  repository, unless prefixed \texttt{crypto-autoresearcher:} or
  \texttt{cryptanalysis:}.
\end{itemize}

\paragraph{What this treatise does not claim.} No speed-up on a deployed
curve, no key recovery, no asymptotic sub-rho algorithm for prime fields,
and no statement about the security of ECC2K-130, P-256 or secp256k1. The
programme's GPU Pollard-rho work on ECC2K-130 is published separately (IACR
ePrint 2026/2469) and is not repeated here.
